% Intended to be inserted near the end of \subsection{Healthcare Resource
% Allocation} (sec:exp-health), after the "Gradient conflict analysis"
% paragraph and before \subsection{Synthetic Multidimensional Knapsack}.
% This paragraph reports the FDFL mu-parameter robustness check and
% keeps the main PCGrad numbers in Table~\ref{tab:healthcare} untouched.

\textit{Sensitivity to the prediction weight $\mu$.}
The scalarized end-to-end family is parameterized by a fixed prediction
weight: the reported FDFL-Scal corresponds to $\mu=1$ in the training
loss $L_{\mathrm{regret}} + \mu\,L_{\mathrm{pred}} + \lambda F$. Two
natural questions are whether $\mu$ needs to be tuned in practice, and
whether the original specification $\mu = 0$ (decision and fairness
only, no explicit MSE anchor) is a viable operating point when the
decision gradient is analytic. To check this, we reran the same
healthcare grid with FDFL at $\mu \in \{0, 0.1, 0.5\}$ alongside the
$\mu = 1$ column already reported. At MAD fairness and $\alpha = 2$,
all four $\mu$ values agree to within $0.003$ on normalized regret
(FDFL $\mu=0$: $0.128 \pm 0.002$; $\mu=0.1$: $0.128 \pm 0.002$;
$\mu=0.5$: $0.129 \pm 0.002$; $\mu=1$: $0.131 \pm 0.002$), and the
same clustering holds for DP, bias-parity, and Atkinson and at
$\alpha = 0.5$. With the analytic decision gradient from
Proposition~\ref{prop:analytical_gradient_group-1}, the regret term
alone already supplies enough signal to keep the predictor calibrated,
so adding an explicit MSE anchor is redundant. Appendix~\ref{app:mu-followup}
documents the full follow-up, including the contrasting behavior in
the synthetic knapsack experiment where $\mu = 0$ diverges under
stochastic decision gradients.
